\section{Biased-SGD descent inequality}
\label{sec:descent_inequality}

\begin{lemma}[Per-step descent inequality]\label{lem:descent_inequality}
Under
\Cref{ass:gradient_correctness_bounded_bias,ass:bounded_second_moment,ass:unembedding_norm},
for every $Q \in F$ and every $j \ge 1$,
\begin{align}\label{eq:descent_inequality}
   \E\!\bigl[\, \loss(x_j; Q) \,\big|\, \mathcal F_{j-1} \,\bigr]
   \;\le\;
   \loss(x_{j-1}; Q)
   \;-\; \tfrac{\eta_j}{2}\, \norm{\nabla \loss(x_{j-1}; Q)}^2
   \;+\; \tfrac{\eta_j}{2}\, \beta^2
   \;+\; \tfrac{L_{\mathrm{sm}}}{2}\, \eta_j^2\, G^2,
\end{align}
where $L_{\mathrm{sm}} \le \tfrac{1}{2} B_U^2$ from
\Cref{lem:smoothness}, and $\beta, \eta_0, \eta_j, G$ are the constants of
\Cref{ass:gradient_correctness_bounded_bias,ass:bounded_second_moment}.
\end{lemma}

\begin{proof}
The proof is the standard biased-SGD descent-inequality argument
(technique digest: \texttt{biased-sgd-descent-inequality.md};
reference: \cite{bottou2018optimization}, \S 4.3, eq.~(4.10)).

\noindent\textbf{Step 1: smoothness expansion.}
By \Cref{lem:smoothness}, $\loss(\cdot; Q)$ is $L_{\mathrm{sm}}$-smooth,
so the second-order Taylor inequality applied to the update
$x_j = x_{j-1} + g_j$ gives, pathwise,
\begin{align}\label{eq:descent_smoothness_step}
   \loss(x_j; Q)
   \;\le\;
   \loss(x_{j-1}; Q)
   \;+\; \inner{\nabla \loss(x_{j-1}; Q)}{g_j}
   \;+\; \tfrac{L_{\mathrm{sm}}}{2}\, \norm{g_j}^2.
\end{align}

\noindent\textbf{Step 2: conditional expectation.}
Take $\E[\,\cdot \mid \mathcal F_{j-1}]$ on both sides of
\Cref{eq:descent_smoothness_step}. Since $x_{j-1}$ (and hence
$\loss(x_{j-1}; Q)$ and $\nabla \loss(x_{j-1}; Q)$) is
$\mathcal F_{j-1}$-measurable, only the inner product and the
$\norm{g_j}^2$ term carry the expectation. Substituting
\Cref{eq:gc_bb_def} into the inner product and
\Cref{eq:second_moment_def} into the squared-norm term,
\begin{align*}
   \E\!\bigl[\, \loss(x_j; Q) \,\big|\, \mathcal F_{j-1} \,\bigr]
   \;\le\;
   \loss(x_{j-1}; Q)
   \;-\; \eta_j \norm{\nabla \loss(x_{j-1}; Q)}^2
   \;+\; \inner{\nabla \loss(x_{j-1}; Q)}{b_j}
   \;+\; \tfrac{L_{\mathrm{sm}}}{2}\, \eta_j^2\, G^2.
\end{align*}

\noindent\textbf{Step 3: Young's inequality on the bias cross-term.}
Apply Young's inequality
$\inner{a}{b} \le \tfrac{\alpha}{2} \norm{a}^2 + \tfrac{1}{2 \alpha} \norm{b}^2$
with $\alpha = \eta_j$ to the inner product
$\inner{\nabla \loss(x_{j-1}; Q)}{b_j}$:
\begin{align*}
   \inner{\nabla \loss(x_{j-1}; Q)}{b_j}
   \;\le\;
   \tfrac{\eta_j}{2} \norm{\nabla \loss(x_{j-1}; Q)}^2
   \;+\; \tfrac{1}{2 \eta_j} \norm{b_j}^2.
\end{align*}
By the relative-scale bias bound \Cref{eq:gc_bb_bias},
$\norm{b_j}^2 \le \eta_j^2 \beta^2$, so
$\tfrac{1}{2 \eta_j} \norm{b_j}^2 \le \tfrac{\eta_j}{2} \beta^2$. Hence
\begin{align*}
   \inner{\nabla \loss(x_{j-1}; Q)}{b_j}
   \;\le\;
   \tfrac{\eta_j}{2} \norm{\nabla \loss(x_{j-1}; Q)}^2
   \;+\; \tfrac{\eta_j}{2}\, \beta^2.
\end{align*}

\noindent\textbf{Step 4: combine.}
Substituting the Young bound into the conditional-expectation display
of Step 2, the gradient-norm terms combine as
$-\eta_j \norm{\nabla \loss}^2 + \tfrac{\eta_j}{2} \norm{\nabla \loss}^2
= -\tfrac{\eta_j}{2} \norm{\nabla \loss}^2$, yielding
\begin{align*}
   \E\!\bigl[\, \loss(x_j; Q) \,\big|\, \mathcal F_{j-1} \,\bigr]
   \;\le\;
   \loss(x_{j-1}; Q)
   \;-\; \tfrac{\eta_j}{2} \norm{\nabla \loss(x_{j-1}; Q)}^2
   \;+\; \tfrac{\eta_j}{2}\, \beta^2
   \;+\; \tfrac{L_{\mathrm{sm}}}{2}\, \eta_j^2\, G^2,
\end{align*}
which is \Cref{eq:descent_inequality}. This completes the proof.
\end{proof}

\begin{remark}\label{rem:descent_inequality_constants}
The descent inequality \Cref{eq:descent_inequality} tracks three
competing terms: the \emph{gradient-descent decrease}
$-\tfrac{\eta_j}{2} \norm{\nabla \loss}^2$, the \emph{bias contribution}
$\tfrac{\eta_j}{2} \beta^2$ (proportional to $\eta_j$), and the
\emph{stochastic noise term} $\tfrac{L_{\mathrm{sm}}}{2} \eta_j^2 G^2$
(proportional to $\eta_j^2$). Both the bias and noise contributions
vanish as $\eta_j \to 0$, but the bias contribution scales as $\eta_j$
while the noise scales as $\eta_j^2$ — under the diminishing-step
schedule of \Cref{lem:telescoping}, the bias term sums to a
$T$-independent floor proportional to $\beta^2$, while the noise term
sums to a $O(1/\sqrt T)$ asymptotic vanishing contribution. The
gradient-descent term drives the loss \emph{down}, and the telescoping
in \Cref{lem:telescoping} extracts the rate at which this descent
dominates the floor.
\end{remark}
